\chapter{Introduction}
\label{ch:introduction}

\section{Motivation}

In 2015 LIGO (Laser Interferometer Gravitational-Wave Observatory) detected
for the first time the collision of two black holes, and in 2017 the gravitational wave
signal GW170817 measured by LIGO and Virgo demonstrated that the fusion of
neutron stars could be observed via telescopes.

Since then, gravitational waves have become a continuous stream of data requiring analysis,
often in real time. In this context, the real challenge is not the faintness of the signals but the noise.
The detectors are so sensitive that they constantly record instrumental and environmental
disturbances known as "glitches".

This is the purpose of Gravity Spy, the source of the data used in this project: to identify
glitch families and trace them back to their causes, as recognizing a glitch
helps in correcting the instrument.

Meanwhile, machine learning has become indispensable for managing data quality,
though it brings with it a cost-related challenge. The demand for computing power
is growing faster than available resources, alerts must be issued within minutes,
and the volume of data is set to increase further.

It is only natural to want to make models more lightweight, and this is where
quantization and spiking neural networks come into play, as they were specifically
designed for low energy consumption.

Deploying a neural network outside a data centre means making it smaller. The
standard tool is indeed quantization, a practice that involves representing weights and activations with fewer
bits, trading numerical precision for memory and energy. The operation is
routinely validated by checking that accuracy does not fall, and for many
applications that is a sufficient test.

Npnetheless, it is not sufficient when the model is used to explain something. Rabbi et
al.~\cite{rabbi2026} evaluated five convolutional architectures at INT8 and
INT4, comparing the attribution maps of the quantized models with those of their
full-precision counterparts. They found that explanation similarity degrades
measurably between INT8 and INT4 even where accuracy differences are negligible,
and that when quantization does cause a misclassification, the quantized model
is generally attending to \emph{different} regions of the input rather than
weighing the same evidence differently. Accuracy and interpretability, in other
words, do not degrade together.

That finding is confined to conventional artificial neural networks. Spiking
neural networks (SNNs), the architecture most closely associated with low-power
deployment~\cite{roy2019towards}, differ from them in ways that make the
question both harder and more interesting.

\section{Research question}

The question this thesis addresses is:

\begin{quote}
\emph{Does a quantized spiking neural network under edge-realistic constraints
base its decisions on the same evidence as its full-precision counterpart, and
at what energy cost?}
\end{quote}

Three properties of spiking networks have no counterpart in the convolutional
setting, and each one bears on the question.

\paragraph{A second quantization target.} A spiking neuron carries a membrane
potential: a state variable read and written at every time step, and a dominant
memory cost at deployment. It is quantized independently of the
weights~\cite{yin2024mint, smith2026}, whereas a conventional network offers only
weights and activations. The two targets can therefore be separated, which turns
a single observation into a causal ablation.

\paragraph{A temporal axis in the explanation.} A spiking explanation is a
\emph{sequence} of maps, one per time step. Quantization may preserve where a
network looks while displacing when it decides, a change to which spatial
metrics are blind by construction.

\paragraph{A native explainer.} The Spike Activation Map~\cite{kim2021sam}
reads spike activity directly. It requires neither gradients, side-stepping the
approximation error a surrogate gradient injects into any transplanted
Grad-CAM-style method~\cite{selvaraju2017gradcam}, nor ground-truth labels.

The working hypothesis follows from the first two properties:

\begin{quote}
\textbf{Quantizing the membrane potential degrades the temporal structure of the
explanation before quantizing the weights degrades its spatial structure},
\end{quote}

since the membrane potential is the variable that integrates evidence over time,
while a weight error is fixed for a whole forward pass. The hypothesis is
falsifiable, and testing it requires quantizing the two targets separately: the
membrane arm should displace the temporal centre of mass of the explanation at a
\emph{higher} bit-width than the one at which the weight arm lowers its spatial
correlation. If the order is reversed, the hypothesis is refuted, which is itself
a reportable result.

\section{Contributions}

Answering the research question requires an instrument, and the metrics used to
compare explanations return numbers between zero and one whatever is fed to them.
Before a quantization result can be read, one must know what those metrics report
when nothing has changed. The thesis is therefore organized in two parts: the
construction and calibration of the instrument, and its application to the
quantization grid.

\begin{enumerate}
  \item \textbf{A hardware-realizable full-precision baseline}, with every
        design choice settled on evidence and recorded in a frozen manifest: the
        input resolution chosen on network health rather than accuracy, the
        membrane decay restricted to shift-implementable values inside the
        hyperparameter search rather than rounded afterwards, and a three-seed
        reference at $0.9833 \pm 0.0014$ validation accuracy.

  \item \textbf{A corrected implementation of the Spike Activation Map}, with
        its decay parameter fixed on the temporal axis it actually governs, and
        its faithfulness tested quantitatively against a random baseline rather
        than judged by visual plausibility.

  \item \textbf{A calibration of the explanation-divergence metrics} at three
        levels: numerical noise, a thirtyfold change in the explainer, and
        initialization variance. Neither of the two closest
        works~\cite{rabbi2026, smith2026} establishes any of the three, and the
        result has the widest scope of the thesis: spatial metrics are valid only
        between a model and its own parent.

  \item \textbf{A bit-width characterization of explanation fidelity for a
        spiking classifier}, with weight and membrane quantization separated and
        then combined, measured with spatial and temporal divergence metrics read
        against the calibrated thresholds, reported per class, and stratified by
        event strength.

  \item \textbf{An analytical energy and memory account} of every operating
        point, so that accuracy, explanation fidelity and energy can be read on
        a common bit-width axis.
\end{enumerate}

\todo{Once the grid is complete, add two or three sentences here stating the
central quantization result (hypothesis confirmed or refuted, and at which
bit-width), mirroring the abstract.}

\section{Application domain}

The classification task is the identification of transient noise --- known as
\emph{glitches} --- in the data of the LIGO gravitational-wave detectors, using
the Gravity Spy release~\cite{zevin2017, glanzer2023, glanzer2021}.

The domain was chosen for three reasons. Explanation fidelity has an
\emph{operational} consequence there: glitch classification guides detector
characterization, where a commissioner acts physically on the instrument, so a
compressed model reaching the same label from different evidence would preserve
the label while corrupting the diagnostic cue. The classes have physically
interpretable morphology, so a domain expert can judge whether an attention map
falls where the glitch actually is, which is not true of a generic image
benchmark. And the data is public, large and labelled.

\section{Scope}

This work characterizes explanation fidelity under \emph{simulated}
quantization, with analytical energy and memory accounting. Hardware deployment
is explicitly out of scope: no model is exported to a neuromorphic device or
microcontroller, and no latency or power is measured on one.

Every design decision in this thesis is taken on the validation partition. The
test partition is opened once, for the final estimates reported in
Section~\ref{sec:test}.

\section{Structure}

Chapter~\ref{ch:background} reviews spiking neural networks, their training by
surrogate gradients, quantization for spiking architectures, explainability
methods applicable to them, and the gravitational-wave classification task.
Chapter~\ref{ch:methodology} describes the data pipeline, the model, the
selection of the operating configuration, the Spike Activation Map and the
protocol used to calibrate the divergence metrics.
Chapter~\ref{ch:quantization} describes the quantization study: the quantizers,
the grid, the fine-tuning protocol, the divergence metrics and the energy model.
Chapter~\ref{ch:results} reports the full-precision reference, the physical
validation of the network dynamics, the explainability analyses, the
calibration, and the results of the quantization grid.
Chapter~\ref{ch:conclusions} discusses the findings and their limits, and
outlines future work.
